% ----------------------------------------------------
%
%        Denotational semantics: Plan A
%
%         Was: desugaring.ltx
%
% ----------------------------------------------------

%% ODER: format ==         = "\mathrel{==}"
%% ODER: format /=         = "\neq "
%
%
\makeatletter
\@ifundefined{lhs2tex.lhs2tex.sty.read}%
  {\@namedef{lhs2tex.lhs2tex.sty.read}{}%
   \newcommand\SkipToFmtEnd{}%
   \newcommand\EndFmtInput{}%
   \long\def\SkipToFmtEnd#1\EndFmtInput{}%
  }\SkipToFmtEnd

\newcommand\ReadOnlyOnce[1]{\@ifundefined{#1}{\@namedef{#1}{}}\SkipToFmtEnd}
\usepackage{amstext}
\usepackage{amssymb}
\usepackage{stmaryrd}
\DeclareFontFamily{OT1}{cmtex}{}
\DeclareFontShape{OT1}{cmtex}{m}{n}
  {<5><6><7><8>cmtex8
   <9>cmtex9
   <10><10.95><12><14.4><17.28><20.74><24.88>cmtex10}{}
\DeclareFontShape{OT1}{cmtex}{m}{it}
  {<-> ssub * cmtt/m/it}{}
\newcommand{\texfamily}{\fontfamily{cmtex}\selectfont}
\DeclareFontShape{OT1}{cmtt}{bx}{n}
  {<5><6><7><8>cmtt8
   <9>cmbtt9
   <10><10.95><12><14.4><17.28><20.74><24.88>cmbtt10}{}
\DeclareFontShape{OT1}{cmtex}{bx}{n}
  {<-> ssub * cmtt/bx/n}{}
\newcommand{\tex}[1]{\text{\texfamily#1}}	% NEU

\newcommand{\Sp}{\hskip.33334em\relax}


\newcommand{\Conid}[1]{\mathit{#1}}
\newcommand{\Varid}[1]{\mathit{#1}}
\newcommand{\anonymous}{\kern0.06em \vbox{\hrule\@width.5em}}
\newcommand{\plus}{\mathbin{+\!\!\!+}}
\newcommand{\bind}{\mathbin{>\!\!\!>\mkern-6.7mu=}}
\newcommand{\rbind}{\mathbin{=\mkern-6.7mu<\!\!\!<}}% suggested by Neil Mitchell
\newcommand{\sequ}{\mathbin{>\!\!\!>}}
\renewcommand{\leq}{\leqslant}
\renewcommand{\geq}{\geqslant}
\usepackage{polytable}

%mathindent has to be defined
\@ifundefined{mathindent}%
  {\newdimen\mathindent\mathindent\leftmargini}%
  {}%

\def\resethooks{%
  \global\let\SaveRestoreHook\empty
  \global\let\ColumnHook\empty}
\newcommand*{\savecolumns}[1][default]%
  {\g@addto@macro\SaveRestoreHook{\savecolumns[#1]}}
\newcommand*{\restorecolumns}[1][default]%
  {\g@addto@macro\SaveRestoreHook{\restorecolumns[#1]}}
\newcommand*{\aligncolumn}[2]%
  {\g@addto@macro\ColumnHook{\column{#1}{#2}}}

\resethooks

\newcommand{\onelinecommentchars}{\quad-{}- }
\newcommand{\commentbeginchars}{\enskip\{-}
\newcommand{\commentendchars}{-\}\enskip}

\newcommand{\visiblecomments}{%
  \let\onelinecomment=\onelinecommentchars
  \let\commentbegin=\commentbeginchars
  \let\commentend=\commentendchars}

\newcommand{\invisiblecomments}{%
  \let\onelinecomment=\empty
  \let\commentbegin=\empty
  \let\commentend=\empty}

\visiblecomments

\newlength{\blanklineskip}
\setlength{\blanklineskip}{0.66084ex}

\newcommand{\hsindent}[1]{\quad}% default is fixed indentation
\let\hspre\empty
\let\hspost\empty
\newcommand{\NB}{\textbf{NB}}
\newcommand{\Todo}[1]{$\langle$\textbf{To do:}~#1$\rangle$}

\EndFmtInput
\makeatother
%
%
%
%
%
%
% This package provides two environments suitable to take the place
% of hscode, called "plainhscode" and "arrayhscode". 
%
% The plain environment surrounds each code block by vertical space,
% and it uses \abovedisplayskip and \belowdisplayskip to get spacing
% similar to formulas. Note that if these dimensions are changed,
% the spacing around displayed math formulas changes as well.
% All code is indented using \leftskip.
%
% Changed 19.08.2004 to reflect changes in colorcode. Should work with
% CodeGroup.sty.
%
\ReadOnlyOnce{polycode.fmt}%
\makeatletter

\newcommand{\hsnewpar}[1]%
  {{\parskip=0pt\parindent=0pt\par\vskip #1\noindent}}

% can be used, for instance, to redefine the code size, by setting the
% command to \small or something alike
\newcommand{\hscodestyle}{}

% The command \sethscode can be used to switch the code formatting
% behaviour by mapping the hscode environment in the subst directive
% to a new LaTeX environment.

\newcommand{\sethscode}[1]%
  {\expandafter\let\expandafter\hscode\csname #1\endcsname
   \expandafter\let\expandafter\endhscode\csname end#1\endcsname}

% "compatibility" mode restores the non-polycode.fmt layout.

\newenvironment{compathscode}%
  {\par\noindent
   \advance\leftskip\mathindent
   \hscodestyle
   \let\\=\@normalcr
   \let\hspre\(\let\hspost\)%
   \pboxed}%
  {\endpboxed\)%
   \par\noindent
   \ignorespacesafterend}

\newcommand{\compaths}{\sethscode{compathscode}}

% "plain" mode is the proposed default.
% It should now work with \centering.
% This required some changes. The old version
% is still available for reference as oldplainhscode.

\newenvironment{plainhscode}%
  {\hsnewpar\abovedisplayskip
   \advance\leftskip\mathindent
   \hscodestyle
   \let\hspre\(\let\hspost\)%
   \pboxed}%
  {\endpboxed%
   \hsnewpar\belowdisplayskip
   \ignorespacesafterend}

\newenvironment{oldplainhscode}%
  {\hsnewpar\abovedisplayskip
   \advance\leftskip\mathindent
   \hscodestyle
   \let\\=\@normalcr
   \(\pboxed}%
  {\endpboxed\)%
   \hsnewpar\belowdisplayskip
   \ignorespacesafterend}

% Here, we make plainhscode the default environment.

\newcommand{\plainhs}{\sethscode{plainhscode}}
\newcommand{\oldplainhs}{\sethscode{oldplainhscode}}
\plainhs

% The arrayhscode is like plain, but makes use of polytable's
% parray environment which disallows page breaks in code blocks.

\newenvironment{arrayhscode}%
  {\hsnewpar\abovedisplayskip
   \advance\leftskip\mathindent
   \hscodestyle
   \let\\=\@normalcr
   \(\parray}%
  {\endparray\)%
   \hsnewpar\belowdisplayskip
   \ignorespacesafterend}

\newcommand{\arrayhs}{\sethscode{arrayhscode}}

% The mathhscode environment also makes use of polytable's parray 
% environment. It is supposed to be used only inside math mode 
% (I used it to typeset the type rules in my thesis).

\newenvironment{mathhscode}%
  {\parray}{\endparray}

\newcommand{\mathhs}{\sethscode{mathhscode}}

% texths is similar to mathhs, but works in text mode.

\newenvironment{texthscode}%
  {\(\parray}{\endparray\)}

\newcommand{\texths}{\sethscode{texthscode}}

% The framed environment places code in a framed box.

\def\codeframewidth{\arrayrulewidth}
\RequirePackage{calc}

\newenvironment{framedhscode}%
  {\parskip=\abovedisplayskip\par\noindent
   \hscodestyle
   \arrayrulewidth=\codeframewidth
   \tabular{@{}|p{\linewidth-2\arraycolsep-2\arrayrulewidth-2pt}|@{}}%
   \hline\framedhslinecorrect\\{-1.5ex}%
   \let\endoflinesave=\\
   \let\\=\@normalcr
   \(\pboxed}%
  {\endpboxed\)%
   \framedhslinecorrect\endoflinesave{.5ex}\hline
   \endtabular
   \parskip=\belowdisplayskip\par\noindent
   \ignorespacesafterend}

\newcommand{\framedhslinecorrect}[2]%
  {#1[#2]}

\newcommand{\framedhs}{\sethscode{framedhscode}}

% The inlinehscode environment is an experimental environment
% that can be used to typeset displayed code inline.

\newenvironment{inlinehscode}%
  {\(\def\column##1##2{}%
   \let\>\undefined\let\<\undefined\let\\\undefined
   \newcommand\>[1][]{}\newcommand\<[1][]{}\newcommand\\[1][]{}%
   \def\fromto##1##2##3{##3}%
   \def\nextline{}}{\) }%

\newcommand{\inlinehs}{\sethscode{inlinehscode}}

% The joincode environment is a separate environment that
% can be used to surround and thereby connect multiple code
% blocks.

\newenvironment{joincode}%
  {\let\orighscode=\hscode
   \let\origendhscode=\endhscode
   \def\endhscode{\def\hscode{\endgroup\def\@currenvir{hscode}\\}\begingroup}
   %\let\SaveRestoreHook=\empty
   %\let\ColumnHook=\empty
   %\let\resethooks=\empty
   \orighscode\def\hscode{\endgroup\def\@currenvir{hscode}}}%
  {\origendhscode
   \global\let\hscode=\orighscode
   \global\let\endhscode=\origendhscode}%

\makeatother
\EndFmtInput
%
%
%
% First, let's redefine the forall, and the dot.
%
%
% This is made in such a way that after a forall, the next
% dot will be printed as a period, otherwise the formatting
% of `comp_` is used. By redefining `comp_`, as suitable
% composition operator can be chosen. Similarly, period_
% is used for the period.
%
\ReadOnlyOnce{forall.fmt}%
\makeatletter

% The HaskellResetHook is a list to which things can
% be added that reset the Haskell state to the beginning.
% This is to recover from states where the hacked intelligence
% is not sufficient.

\let\HaskellResetHook\empty
\newcommand*{\AtHaskellReset}[1]{%
  \g@addto@macro\HaskellResetHook{#1}}
\newcommand*{\HaskellReset}{\HaskellResetHook}

\global\let\hsforallread\empty

\newcommand\hsforall{\global\let\hsdot=\hsperiodonce}
\newcommand*\hsperiodonce[2]{#2\global\let\hsdot=\hscompose}
\newcommand*\hscompose[2]{#1}

\AtHaskellReset{\global\let\hsdot=\hscompose}

% In the beginning, we should reset Haskell once.
\HaskellReset

\makeatother
\EndFmtInput
% ----------------------
%% Stuff for TimCore


%% 'g' inferences





%% Intersection of effects

%% --------- Effects ------------


%% ----------------------


%% --------------------------
%%        Desugaring rules
%% --------------------------



% ----------------------------------------------------
% New desugaring, just does the wrapping part:
% ----------------------------------------------------
%   Input for wrapping


%%format Hide (e) = "{\mathcal H}\llbracket{}" e "\rrbracket{}"

%% 'si' is short for "solve or infer" or "unification variable or not" resp

%% 'md' is short for "mode" or "unification variable or not" resp

%% flipsi is a no-op for now


%% EXX is just EX, but with no argument




%% ----------- Applications --------------
%% %format applyx a b = a "@" b
%% %format applyx a b = a "\," b

%% Function application with no arguments

%% Function application with 1 argument

%% Function application with 1 argument, no parens

%% Function application with 1 argument, without parenthesis

%% Function application with many arguments

%% ----------- End of Applications --------------





%% Sequences e1; ...; en; v

%% xs etc stole syntax used in examples.

%% Old version: | for BNF, [] for choice
%%%format choice = "[ \! ]"
%%%format `choice` = "\mathop{[ \! ]}"
%% 1mm is about 0.3em with the current font

%% New version: tall | for BNF, fat | for choice
%%format vbar = "\mathop{\bigm|}"
%%format choice = "\boldsymbol{|}"
%%format `choice` = "\mathop{\boldsymbol{|}}"




%% %format ok = "\textbf{ok}"

%% %format `eqsemi` = "\hspace{-0.5mm}\fatsemi{}"
%% %format eqsemiKW = "\fatsemi"

%% Arrays and tuples

%% %format vmap (x) = "\llangle" x "\rrangle"






%%  OLD SYNTAX   %format split (x) (y) (z) = "\textbf{split}(" x ")\{" y "," z "\}"
%% %format defKW = "\textbf{def}"
% only used in comments, but still prettier this way:
%%%%format map = "\textbf{map}"

%% -------Functions ---------------


%% functionOC has open/closed subscript


%% -------End of f unctions ---------------





%% Unification
%%%format == = "\!=\!"


% format := = "\mapsto"
% %format squiggt = "\leadsto"
% %format squiggt = "\!-\!\!>\!\!\!"



































%% We combined S and C into one


%% ----------- Metatheory --------------


%% ----------------------------
%%         Verification
%% ----------------------------

%% assert = succeed, abbreviated to suc



%% Effects checking, like succeeds{e}





%% <fx> brackets

\begin{figure}
$$
% ------------- Syntax of Mini Verse ---------------
\begin{array}{l}
\textbf{Syntax of Mini Verse} \\
\begin{array}{lr@{\;\;}c@{\;\;}l}
  \text{Expressions} & e & ::= & \ensuremath{\Varid{x}} \vb \ensuremath{\Varid{k}} \vb \ensuremath{\Varid{o}} \vb \ensuremath{\Varid{e}_{\mathrm{1}} [\!\Varid{e}_{\mathrm{2}}]} \vb e_1; e_2
     \vb \ensuremath{\Varid{e}_{\mathrm{1}}\hspace{-0.14em}=\hspace{-0.14em}\Varid{e}_{\mathrm{2}}} \vb \ensuremath{\exists\Varid{x}} \vb \ensuremath{\Varid{e}_{\mathrm{1}}\hspace{0.19em}\mathop{\rule[-0.1em]{0.15em}{0.75em}}\hspace{0.2em}\Varid{e}_{\mathrm{2}}} \vb \ensuremath{\textbf{fail}}
     \vb \ensuremath{\lambda_{\Varid{q}}\Varid{i}.(\Varid{e}_{\mathrm{1}})\{\Varid{e}_{\mathrm{2}}\}} \\
  && \vb & \ensuremath{\langle\Varid{e}_{\mathrm{1}},\mydots,\Varid{e}_{\Varid{n}}\rangle} \vb \ensuremath{\textbf{truth}\{\Varid{e}\}} \vb \ensuremath{\Varid{t}_{\mathrm{1}}\vartriangleright\hspace{-1.3mm} \langle\omega\rangle\,\Varid{t}_{\mathrm{2}}} \vb \ensuremath{{\bf chk}\langle\omega\rangle\{\Varid{e}\}} \\
  && \vb & \ensuremath{\textbf{if}(\Varid{t}_{\mathrm{1}})\{\Varid{t}_{\mathrm{2}}\}\textbf{else}\{\Varid{t}_{\mathrm{3}}\}} \vb \ensuremath{\textbf{for}(\Varid{t}_{\mathrm{1}})\,\textbf{do}\,\Varid{t}_{\mathrm{2}}} \vb \ensuremath{\textbf{all}\{\Varid{e}\}} \\
\multicolumn{4}{l}{\hspace{3mm}\text{Note: no \ensuremath{\mathbin{:}\Varid{t}}; no \ensuremath{\Varid{x}\!\rightarrow\!\Varid{t}}; no \ensuremath{\Varid{x}\mathrel{\coloneqq}\Varid{e}}; no \ensuremath{\mathbf{where}}; and lambda changes form}}
\end{array} \\
\\
% ------------- Desugaring ---------------
\begin{array}{@{}l@{\,}r@{\;\;}c@{\;\;}l@{}l}
  \text{Inputs} & \ensuremath{\kappa} & ::= & \ensuremath{\Varid{dr},\!\Varid{u}} \\
  & \ensuremath{\Varid{dr}} & ::= & \multicolumn{2}{l}{\ensuremath{D}  \vb \ensuremath{R\langle\omega\rangle}
     \hspace{1cm} \text{\ensuremath{D} for domain, \ensuremath{R\langle\omega\rangle} for range}} \\
                & \ensuremath{\Varid{u}}  & ::= & \ensuremath{\bullet} \vb \ensuremath{\Varid{e}} \\[2mm]

\rulename{wconst}      & \ensuremath{{\mathcal W}\llbracket\Varid{k}\rrbracket\kappa} & = & \ensuremath{\kappa\;\mathring{=}\;\Varid{k}} \\
\rulename{wprim}   & \ensuremath{{\mathcal W}\llbracket\Varid{o}\rrbracket\kappa} & = & \ensuremath{\kappa\;\mathring{=}\;\Varid{o}} \\
\rulename{wvar}    & \ensuremath{{\mathcal W}\llbracket\Varid{x}\rrbracket\kappa} & = & \ensuremath{\kappa\;\mathring{=}\;\Varid{x}} \\
\rulename{wwild}    & \ensuremath{{\mathcal W}\llbracket\anonymous \rrbracket\kappa} & = & \ensuremath{\kappa\;\mathring{=}\;(\exists\Varid{x};\,\Varid{x})} \\
% \rulename{wwild2}    & |SS (apeff e fx) _| & = & |e| \\
\rulename{wapp}    & \ensuremath{{\mathcal W}\llbracket\Varid{t}_{\mathrm{1}} [\!\Varid{t}_{\mathrm{2}}]\rrbracket\kappa} & = & \ensuremath{\kappa\;\mathring{=}\;({\mathcal W}\llbracket\Varid{t}_{\mathrm{1}}\rrbracket(D,\bullet)) [\!{\mathcal W}\llbracket\Varid{t}_{\mathrm{2}}\rrbracket(D,\bullet)]} \\
\rulename{wfor}   & \ensuremath{{\mathcal W}\llbracket\textbf{for}(\Varid{t}_{\mathrm{1}})\,\textbf{do}\,\Varid{t}_{\mathrm{2}}\rrbracket\kappa}
   & = & \ensuremath{\kappa\;\mathring{=}\;\textbf{for}({\mathcal W}\llbracket\Varid{t}_{\mathrm{1}}\rrbracket(D,\bullet))\,\textbf{do}\,{\mathcal W}\llbracket\Varid{t}_{\mathrm{2}}\rrbracket(D,\bullet)} \\
\rulename{wone} & \ensuremath{{\mathcal W}\llbracket\textbf{one}\{\Varid{t}\}\rrbracket\kappa} & = & \ensuremath{\kappa\;\mathring{=}\;\textbf{one}\{{\mathcal W}\llbracket\Varid{t}\rrbracket(D,\bullet)\}} \\
\rulename{wall} & \ensuremath{{\mathcal W}\llbracket\textbf{all}\{\Varid{t}\}\rrbracket\kappa} & = & \ensuremath{\kappa\;\mathring{=}\;\textbf{all}\{{\mathcal W}\llbracket\Varid{t}\rrbracket(D,\bullet)\}} \\
\rulename{wfail}   & \ensuremath{{\mathcal W}\llbracket\textbf{fail}\rrbracket\kappa} & = & \ensuremath{\textbf{fail}} \\

% Choice: do we want to push wu into the branches or not?
% \rulename{wchoice} \asktim & |SS wu (t_1 `choice` t_2)| & = & |wu `ueq` (SS wuinit t_1 `choice` SS wuinit t_2)| \\
\rulename{wchoice} \asktim & \ensuremath{{\mathcal W}\llbracket\Varid{t}_{\mathrm{1}}\hspace{0.19em}\mathop{\rule[-0.1em]{0.15em}{0.75em}}\hspace{0.2em}\Varid{t}_{\mathrm{2}}\rrbracket\kappa} & = & \ensuremath{{\mathcal W}\llbracket\Varid{t}_{\mathrm{1}}\rrbracket\kappa\hspace{0.19em}\mathop{\rule[-0.1em]{0.15em}{0.75em}}\hspace{0.2em}{\mathcal W}\llbracket\Varid{t}_{\mathrm{2}}\rrbracket\kappa} \\

\rulename{wunify}   & \ensuremath{{\mathcal W}\llbracket\Varid{t}_{\mathrm{1}}\hspace{-0.14em}=\hspace{-0.14em}\Varid{t}_{\mathrm{2}}\rrbracket\kappa} & = & \ensuremath{{\mathcal W}\llbracket\Varid{t}_{\mathrm{1}}\rrbracket\kappa\mathrel{=}{\mathcal W}\llbracket\Varid{t}_{\mathrm{2}}\rrbracket\kappa} \\
\rulename{wsemi}\asktim   & \ensuremath{{\mathcal W}\llbracket\Varid{t}_{\mathrm{1}};\,\Varid{t}_{\mathrm{2}}\rrbracket\kappa} & = & \ensuremath{{\mathcal W}\llbracket\Varid{t}_{\mathrm{1}}\rrbracket(D,\bullet);\,{\mathcal W}\llbracket\Varid{t}_{\mathrm{2}}\rrbracket\kappa} \\
\rulename{wwhere}   & \ensuremath{{\mathcal W}\llbracket\Varid{t}_{\mathrm{1}}\;\mathbf{where}\;\Varid{t}_{\mathrm{2}}\rrbracket\kappa} & = & \ensuremath{\exists\Varid{r};\,\Varid{r}\hspace{-0.14em}=\hspace{-0.14em}{\mathcal W}\llbracket\Varid{t}_{\mathrm{1}}\rrbracket\kappa;\,{\mathcal W}\llbracket\Varid{t}_{\mathrm{2}}\rrbracket(D,\bullet);\,\Varid{r}} \\
\multicolumn{2}{@{}l}{\rulename{wif} \; \hfill \ensuremath{{\mathcal W}\llbracket\textbf{if}(\Varid{t}_{\mathrm{1}})\{\Varid{t}_{\mathrm{2}}\}\textbf{else}\{\Varid{t}_{\mathrm{3}}\}\rrbracket\kappa}}
 & = & \ensuremath{\textbf{if}({\mathcal W}\llbracket\Varid{t}_{\mathrm{1}}\rrbracket(D,\bullet))\{{\mathcal W}\llbracket\Varid{t}_{\mathrm{2}}\rrbracket\kappa\}\textbf{else}\{{\mathcal W}\llbracket\Varid{t}_{\mathrm{3}}\rrbracket\kappa\}} \\
\rulename{wdef}   & \ensuremath{{\mathcal W}\llbracket\Varid{x}\mathrel{\coloneqq}\Varid{t}\rrbracket\kappa} & = & \ensuremath{\exists\Varid{x};\,\Varid{x}\mathrel{=}{\mathcal W}\llbracket\Varid{t}\rrbracket\kappa} \\
\rulename{wsquig}   & \ensuremath{{\mathcal W}\llbracket\Varid{x}\!\rightarrow\!\Varid{t}\rrbracket(\Varid{dr},\!\Varid{u})}
  & = & \ensuremath{(\Varid{dr},\!\Varid{u})\;\mathring{=}\;\Varid{x};\,{\mathcal W}\llbracket\Varid{t}\rrbracket(\Varid{dr},\!\Varid{x})} \\
% \rulename{wsquig2}   & |SS (apeff e fx) (squig x y := t)| & = & |defintro (x y); x == e; y == SS (apeff x fx) t| \\
\rulename{wchk}   & \ensuremath{{\mathcal W}\llbracket{\bf chk}\langle\omega\rangle\{\Varid{t}\}\rrbracket\kappa} & = & \ensuremath{{\bf chk}\langle\omega\rangle\{{\mathcal W}\llbracket\Varid{t}\rrbracket\kappa\}} \\
\rulename{wcol1}\asktim   & \ensuremath{{\mathcal W}\llbracket\mathbin{:}\Varid{t}\rrbracket(\Varid{dr},\!\bullet)} & = & \ensuremath{\exists\Varid{x};\,({\mathcal W}\llbracket\Varid{t}\rrbracket(D,\bullet)) [\!\Varid{x}]} & \text{\ensuremath{\Varid{x}} fresh}\\
\rulename{wcol2}   & \ensuremath{{\mathcal W}\llbracket\mathbin{:}\Varid{t}\rrbracket(D,\!\Varid{e})} & = & \ensuremath{({\mathcal W}\llbracket\Varid{t}\rrbracket(D,\bullet)) [\!\Varid{e}]} \\
\rulename{wcol3}   & \ensuremath{{\mathcal W}\llbracket\mathbin{:}\Varid{t}\rrbracket(R\langle\omega\rangle,\!\Varid{e})} & = & \ensuremath{\Varid{e}\vartriangleright\hspace{-1.3mm} \langle\omega\rangle\,{\mathcal W}\llbracket\Varid{t}\rrbracket(D,\bullet)} \\
\rulename{wopaque}\asktim   & \ensuremath{{\mathcal W}\llbracket\Varid{t}_{\mathrm{1}}\vartriangleright\hspace{-1.3mm} \langle\omega_{\mathrm{1}}\rangle\,\Varid{t}_{\mathrm{2}}\rrbracket(R\langle\omega_{\mathrm{2}}\rangle,\!\Varid{u})}
    & = & \ensuremath{{\mathcal W}\llbracket\Varid{t}_{\mathrm{1}}\rrbracket(D,\!\Varid{u})\vartriangleright\hspace{-1.3mm} \langle\omega_{\mathrm{1}}\cap\omega_{\mathrm{2}}\rangle\,{\mathcal W}\llbracket\Varid{t}_{\mathrm{2}}\rrbracket(D,\bullet)} \\
\rulename{wtup1}   & \ensuremath{{\mathcal W}\llbracket\langle\Varid{t}_{\mathrm{1}},\mydots,\Varid{t}_{\Varid{n}}\rangle\rrbracket(\Varid{dr},\!\bullet)}
   & = & \ensuremath{\langle{\mathcal W}\llbracket\Varid{t}_{\mathrm{1}}\rrbracket(\Varid{dr},\!\bullet),\mydots,{\mathcal W}\llbracket\Varid{t}_{\Varid{n}}\rrbracket(\Varid{dr},\!\bullet)\rangle} \\
\rulename{wtup2}   & \ensuremath{{\mathcal W}\llbracket\langle\Varid{t}_{\mathrm{1}},\mydots,\Varid{t}_{\Varid{n}}\rangle\rrbracket(\Varid{dr},\!\Varid{e})}
& = & \begin{array}[t]{@{}l}
  \ensuremath{\exists\Varid{x}_{\mathrm{1}}\;\mydots\;\Varid{x}_{\Varid{n}};\,\Varid{e}\mathrel{=}\langle\Varid{x}_{\mathrm{1}},\mydots,\Varid{x}_{\Varid{n}}\rangle;\,} \\
  \ensuremath{\langle{\mathcal W}\llbracket\Varid{t}_{\mathrm{1}}\rrbracket(\Varid{dr},\!\Varid{x}_{\mathrm{1}}),\mydots,{\mathcal W}\llbracket\Varid{t}_{\Varid{n}}\rrbracket(\Varid{dr},\!\Varid{x}_{\Varid{n}})\rangle}
  \end{array} & \text{\ensuremath{\Varid{x}_{\Varid{i}}} fresh}\\
\rulename{wtru1} & \ensuremath{{\mathcal W}\llbracket\textbf{truth}\{\Varid{t}\}\rrbracket(\Varid{dr},\!\bullet)}
   & = & \ensuremath{\textbf{truth}\{{\mathcal W}\llbracket\Varid{t}\rrbracket(\Varid{dr},\!\bullet)\}} \\
\rulename{wtru2} & \ensuremath{{\mathcal W}\llbracket\textbf{truth}\{\Varid{t}\}\rrbracket(\Varid{dr},\!\Varid{e})}
   & = & \ensuremath{\exists\Varid{x};\,\Varid{e}\hspace{-0.14em}=\hspace{-0.14em}\textbf{truth}\{\Varid{x}\};\,\textbf{truth}\{{\mathcal W}\llbracket\Varid{t}\rrbracket(\Varid{dr},\!\Varid{x})\}} & \text{\ensuremath{\Varid{x}} fresh} \\
\rulename{wfun1} \asktim   & \ensuremath{{\mathcal W}\llbracket\textbf{fun}_{\Varid{q}}(\Varid{t}_{\mathrm{1}})\langle\omega\rangle\{\Varid{t}_{\mathrm{2}}\}\rrbracket(\Varid{dr},\!\bullet)}
   & = & \ensuremath{\lambda_{\Varid{q}}\Varid{i}.({\mathcal W}\llbracket\Varid{t}_{\mathrm{1}}\rrbracket(D,\!\Varid{i}))\{{\bf chk}\langle\omega\rangle\{{\mathcal W}\llbracket\Varid{t}_{\mathrm{2}}\rrbracket(R\langle\omega\rangle,\!\bullet)\}\}} & \text{\ensuremath{\Varid{i}} fresh} \\
\rulename{wfun2}   & \ensuremath{{\mathcal W}\llbracket\textbf{fun}_{\Varid{q}}(\Varid{t}_{\mathrm{1}})\langle\omega\rangle\{\Varid{t}_{\mathrm{2}}\}\rrbracket(\Varid{dr},\!\Varid{f})}
& = & \begin{array}[t]{@{}l}
  \ensuremath{\textbf{isFun} [\!\Varid{f}];\,} \\
  \ensuremath{\lambda_{\Varid{q}}\Varid{i}.\begin{array}[t]{@{}l} (\exists\Varid{x};\,\Varid{x}\mathrel{=}{\mathcal W}\llbracket\Varid{t}_{\mathrm{1}}\rrbracket(D,\!\Varid{i})) \\ \{{\mathcal K}_{\Varid{q}}\llbracket\Varid{t}_{\mathrm{1}}\rrbracket\Varid{f};\,{\bf chk}\langle\omega\rangle\{{\mathcal W}\llbracket\Varid{t}_{\mathrm{2}}\rrbracket(R\langle\omega\rangle,\!\Varid{f} [\!\Varid{x}])\}\} \end{array}}
  \end{array} & \text{\ensuremath{\Varid{i},\Varid{x}} fresh}
\end{array} \\[1mm]
\textbf{Auxiliary functions} \\
\hspace{3cm}
\begin{array}{@{}l@{\;\;}r@{\;\;}c@{\;\;}ll}
\rulename{wp1} & \ensuremath{(\Varid{dr},\!\bullet)\;\mathring{=}\;\Varid{e}_{\mathrm{2}}} & = & \ensuremath{\Varid{e}_{\mathrm{2}}} \\
\rulename{wp2} \asktim & \ensuremath{(D,\!\Varid{e}_{\mathrm{1}})\;\mathring{=}\;\Varid{e}_{\mathrm{2}}} & = & \ensuremath{\Varid{e}_{\mathrm{1}}\hspace{-0.14em}=\hspace{-0.14em}\Varid{e}_{\mathrm{2}}} \\
\rulename{wp3} \asktim & \ensuremath{(R\langle\omega\rangle,\!\Varid{e}_{\mathrm{1}})\;\mathring{=}\;\Varid{e}_{\mathrm{2}}}
 & = & \ensuremath{\Varid{e}_{\mathrm{1}}\vartriangleright\hspace{-1.3mm} \langle \rangle\,(\lambda\Varid{i} .\,\Varid{i}\hspace{-0.14em}=\hspace{-0.14em}\Varid{e}_{\mathrm{2}})} & \text{$\ensuremath{\Varid{i}} \not\in \ensuremath{\mathsf{fvs}(\Varid{e}_{\mathrm{2}})}$}
\\[2mm]
& \ensuremath{{\mathcal K}_{\textbf{o}}\llbracket\Varid{e}\rrbracket\Varid{f}} & = & \ensuremath{\langle\rangle} \\
& \ensuremath{{\mathcal K}_{\textbf{c}}\llbracket\Varid{e}\rrbracket\Varid{f}} & = & \ensuremath{\lambda_{\textbf{c}}\Varid{x}.(\Varid{f} [\!\Varid{x}])\{\Varid{x}\hspace{-0.14em}=\hspace{-0.14em}{\mathcal W}\llbracket\Varid{e}\rrbracket(D,\bullet)\}}
\end{array}
\end{array}
$$
\caption{Plan C: Desugaring Essential Verse to Mini Verse} \label{fig:wrapping}
\end{figure}

Key points about \Cref{fig:wrapping}, the Essential-Verse-to-Mini-Verse wrapping step:
\begin{itemize}
\item Essential Verse has \ensuremath{(\Varid{t}_{\mathrm{1}}\;\mathbf{where}\;\Varid{t}_{\mathrm{2}})} as well as \ensuremath{(\Varid{t}_{\mathrm{1}};\,\Varid{t}_{\mathrm{2}})}.  We can't flip the
  order without risking getting choices in the wrong order.  But aft
er the wrapping
  step we can replace \ensuremath{\mathbf{where}} by ordering ``\ensuremath{;\,}'', because we have by then
  feed \ensuremath{\Varid{u}} into \ensuremath{\Varid{t}_{\mathrm{1}}} rather than \ensuremath{\Varid{t}_{\mathrm{2}}}, which is the key difference.
\end{itemize}

